\documentclass[../../../include/open-logic-section]{subfiles}

\begin{document}

\olfileid[pt]{sth}{ordinals}{vn}	
\olsection[Construção de Von Neumann]{Construção de Von Neumann dos Ordinais}

\olref[sth][ordinals][iso]{thm:woalwayscomparable} sugere uma ideia.
Poderíamos introduzir certos objetos, chamados \emph{tipos de ordem},
para servirem de representantes das boas ordenações. Escrevendo
$\ordtype{A, <}$ para o tipo de ordem da boa ordenação $\tuple{A, <}$,
gostaríamos de garantir os dois princípios seguintes:
\begin{align*}
	\ordtype{A, <} = \ordtype{B, \lessdot} & 
	\text{ sse } \ordeq{\tuple{A, <}}{\tuple{B, \lessdot}}\\
	\ordtype{A, <} < \ordtype{B, \lessdot}&
	\text{ sse }\ordeq{\tuple{A, <}}{\tuple{B_b, \lessdot_b}}\text{ para algum }b \in B
\end{align*}
Além disso, poderíamos esperar introduzir tipos de ordem \emph{como
certos conjuntos}, tal como podemos introduzir os números naturais como
certos conjuntos.

A forma mais comum de o fazer---e a abordagem que seguiremos---é definir
estes tipos de ordem via certos conjuntos bem ordenados
\emph{canónicos}. Estes conjuntos canónicos foram introduzidos pela
primeira vez por von Neumann:

\begin{defn}
O conjunto $A$ é \emph{transitivo} {sse} $(\forall x \in A)x \subseteq
A$. Então $A$ é um \emph{ordinal} {sse} $A$ é transitivo e bem ordenado
por $\in$.
\end{defn}
\noindent
Daqui em diante, usaremos letras gregas para ordinais. Segue-se
imediatamente da definição que, se $\alpha$ é um ordinal, então
$\tuple{\alpha, \in_\alpha}$ é uma boa ordenação, onde $\in_\alpha =
\Setabs{\tuple{x, y} \in \alpha^2}{x \in y}$. Assim, abusando um pouco da
notação, podemos dizer que o próprio $\alpha$ \emph{é} uma boa ordenação.

Eis os primeiros ordinais:
\[
	\emptyset, \{\emptyset\}, 
	\{\emptyset, \{\emptyset\}\}, 
	\{\emptyset, \{\emptyset\}, \{\emptyset, \{\emptyset\}\}\}, \ldots
\]
Repare que estes são precisamente os primeiros ordinais que
encontramos no nosso Axioma do Infinito, isto é, na definição de von
Neumann de $\omega$ (ver \olref[sth][z][infinity-again]{sec}). Isto não é
coincidência. A definição de von Neumann dos ordinais trata os números
naturais como ordinais, mas permite também ordinais transfinitos.

Como sempre, podemos agora perguntar: \emph{são} estes os ordinais? Ou
von Neumann deu-nos apenas alguns conjuntos que podemos \emph{tratar}
como os ordinais? O tipo de discussões que se podem ter sobre esta
questão é semelhante às que tivemos na \olref[sfr][rel][ref]{sec},
\olref[sfr][arith][ref]{sec},
\olref[sfr][infinite][dedekindsproof]{sec}, e
\olref[sth][z][nat]{sec}, pelo que não insistiremos no ponto.
Em vez disso, daqui em diante usaremos ``os ordinais'' para falar dos
``ordinais de von Neumann''.

\end{document}
